% META: {"id":"1SPE-TRIGO-CO-005","chapitre":"1SPE-TRIGONOMETRIE","type_objet":"corrige","exercice_ref":"1SPE-TRIGO-EX-005","capacites_codes":["C1"],"sources_inspiration":[],"status":"approved","origin":"generated","status_history":["generated","audited","accepted","approved"],"release_acceptance":"RELEASE_OWNER_BATCH_ACCEPTANCE","acceptance_closure_digest":"sha256:9b3ccf9a81c5520fbb7e03b7b2d3e7bf2057a02834b6d908bf3be5f49f3b553f","acceptance_receipt_digest":"sha256:4fad86f0282e0fa4e0419cca1ad6379ae7af5a280c04a4a90880f90a1c044242"}
% BEGIN-VERIFY
% from sympy import *
% assert Rational(100,3) - 16*2 == Rational(4,3)
% assert Rational(4,3) - 2 == Rational(-2,3)
% assert Rational(-23,4) + 6 == Rational(1,4)
% assert 150 * pi / 180 == 5*pi/6
% END-VERIFY
\begin{corrige}{1SPE-TRIGO-EX-005}

\textbf{1.} $\dfrac{100\pi}{3}$ : on calcule $\dfrac{100}{3 \times 2} = \dfrac{100}{6} \approx 16{,}67$. Donc on retranche $16 \times 2\pi$ :
\[
\dfrac{100\pi}{3} - 32\pi = \dfrac{100\pi - 96\pi}{3} = \dfrac{4\pi}{3}.
\]
Or $\dfrac{4\pi}{3} > \pi$, donc $\dfrac{4\pi}{3} - 2\pi = -\dfrac{2\pi}{3}$.

La mesure principale est $-\dfrac{2\pi}{3}$.

\medskip
\textbf{2.} $-\dfrac{23\pi}{4}$ : on calcule $\dfrac{23}{4 \times 2} = \dfrac{23}{8} \approx 2{,}875$. On ajoute $3 \times 2\pi$ :
\[
-\dfrac{23\pi}{4} + 6\pi = -\dfrac{23\pi}{4} + \dfrac{24\pi}{4} = \dfrac{\pi}{4}.
\]

La mesure principale est $\dfrac{\pi}{4}$.

\medskip
\textbf{3.} $150° = 150 \times \dfrac{\pi}{180} = \dfrac{5\pi}{6}$.

\medskip
\textbf{4.} Le point $(-1;0)$ correspond a l'angle $\pi$. La mesure principale est $\pi$.

\end{corrige}
